\section{Further arithmetic reductions}

\paragraph{Scope and provenance.}
The source is Version~10 at commit\\
{\small\nolinkurl{036cebe4f1a4aa7ca89820a9b48f3629b4c364d6}}.\\
The starting system for this note is the previously checked construction
with the packed parity anchor. We describe four further changes:
six field-valued difference guards become rational guards; the scalar
correction becomes the fixed integer $21$; the unused scalar factor
$2\beta-1$ is removed; and the scalar preparation uses a conic instead
of its rational parametrization.
The guard and preparation changes have explicit two-way witness maps on
the stated domains. The input-label and correction change instead reuses
the source's global soundness and completeness argument with a new label.
It retains every inherited arithmetic dependency; it is not an independent
proof of the source's global characterization.

\subsection{Six rational guards for elliptic differences}

Put $\rho=\sqrt2$, $K=\mathbb Q(\rho)$, $D=11-6\rho$, and
\[
 E:y^2=x^3-Dx,\qquad
 P=\left(\frac{57}{8}-\frac72\rho,
          -\frac{165}{8}+\frac{417}{32}\rho\right).
\]
The inherited point-domain result identifies $2E(K)=\mathbb ZP$ and
makes every chart point a finite nonzero multiple of $P$.
We use $z=-x/y$ and normalize the valuation of $K$ above $2$ by $v(2)=1$.
For $a,b\in\mathbb Q$, the nonzero terms in $a+b\rho$ have valuations
$v_2(a)$ and $v_2(b)+\tfrac12$. They have different fractional parts, so
\begin{equation}\label{further:eq:componentvaluation}
 v(a+b\rho)=\min\{v_2(a),v_2(b)+\tfrac12\}.
\end{equation}

The chart values
\[
 h(P)=-1+\frac74\rho,\qquad
 j(P)=\frac{39}{8}-\frac92\rho,\qquad z(P)=-h(P)/j(P)
\]
give $v(z(P))=\tfrac32$. The formal group is defined over the local
integer ring. Its inverse is $T\mapsto-T$, and
$[u](T)=uT+O(T^3)$. Odd multiplication preserves a depth at least
$\tfrac32$. Multiplication by $2$ increases it by exactly one, since
$3r>1+r$ at such a depth $r$. Consequently, for every nonzero integer $n$,
\begin{equation}\label{further:eq:multipledepth}
 t_n=z(nP),\qquad v(t_n)=\frac32+v_2(n).
\end{equation}
Negative integers are included because $t_{-n}=-t_n$.

Write the formal group law as $F(X,Y)$. The identity $F(X,-X)=0$ and
its linear term give
\[
 F(X,Y)=(X+Y)A(X,Y),\qquad
 A\in\mathcal O_K[[X,Y]],\quad A(0,0)=1.
\]
Thus, when $n\ne\pm m$, both evaluated factors $A(t_n,\pm t_m)$ are
units and
\[
 v(t_n-t_m)=\frac32+v_2(n-m),\qquad
 v(t_n+t_m)=\frac32+v_2(n+m).
\]
The integral Laurent expansion
$x(T)=T^{-2}+\sum_{j\ge1}a_jT^{2j}$ likewise yields
\[
 x(t)-x(s)=\frac{s^2-t^2}{s^2t^2}\,B(t,s),
 \qquad B\in\mathcal O_K[[t,s]],\quad B(0,0)=1.
\]
It follows that
\begin{equation}\label{further:eq:differencevaluation}
 \begin{split}
 v\bigl(x(nP)-x(mP)\bigr)
 ={}&-3+v_2(n-m)+v_2(n+m)\\
    &-2v_2(n)-2v_2(m)\in\mathbb Z.
 \end{split}
\end{equation}
If the nonzero difference is $a+b\rho$, equations
\eqref{further:eq:componentvaluation} and
\eqref{further:eq:differencevaluation} force $a\ne0$.

Therefore, on this inherited point domain,
\[
 x(Q)\ne x(R)
 \quad\Longleftrightarrow\quad
 \operatorname{coeff}_{1}\bigl(x(Q)-x(R)\bigr)\ne0.
\]
Replace each of the six unused field inverses of a subtraction denominator
by one rational variable $g$ satisfying $ag=1$.
The two slope coordinates are retained. An old solution determines
$g=1/a$; a new solution has a nonzero field denominator and recovers its
old inverse uniquely. The cases $n=\pm m$ have zero denominator and are
excluded by both guards. This saves six coordinates without altering the
projected points or any quotient value.

\subsection{A local anchor that also checks the input at 2 and 3}

Let the four root points be $Q_i=[n_i]P$, $n_i\ne0$, and write
\[
 \prod_{i=1}^{4}x(Q_i)=C_0+C_1\rho,
 \qquad B_0=4096C_0.
\]
Let $S_{\rm pt}=[a]P$, $R=S_{\rm pt}-P$, and
$f(R)=1/x(R)=f_0+f_1\rho$. Set
\[
 C_{\rm a}=3^{2^{25}},\qquad
 \mathsf S=C_{\rm a}^{2}=3^{2^{26}},\qquad
 q_2(u,v)=u^2+uv-v^2.
\]
The point subtraction guard ensures $R$ is finite and nonidentity.

For completeness, the local facts behind the earlier packed anchor are
as follows. At $3$, the extension $K/\mathbb Q$ is unramified and inert,
with residue field $\mathbb F_9$. The reduction of $P$ is $(\rho,0)$,
of order two, and direct doubling gives
\[
 v_3(z(2P))=1,\qquad z(2P)/3\equiv2+\rho\pmod3.
\]
For an even nonzero $n=2\cdot3^ku$, $3\nmid u$, the integral odd
formal group law gives
\[
 3^{-(k+1)}z(nP)\equiv u(2+\rho)\pmod3.
\]
Indeed, multiplication by a $3$-adic unit scales the normalized residue
by that unit, and multiplication by $3$ raises depth by one without
changing the normalized residue. The error after division by the new
power of $3$ has valuation at least one. These assertions also hold for
negative $u$.
Since $xz^2\equiv1$ on the formal branch,
$u^2\equiv1\pmod3$, and $(2+\rho)^{-2}=2\rho$ in $\mathbb F_9$,
\[
 3^{e(n)}x(nP)\equiv
 \begin{cases}
 \rho,&n\text{ odd},\\
 2\rho,&n\text{ even},
 \end{cases}
 \quad
 e(n)=\begin{cases}0,&n\text{ odd},\\2+2v_3(n),&n\text{ even}.
 \end{cases}
\]
Writing $E_0=\sum_i e(n_i)$ and letting $j$ count the even indices gives
\[
 3^{E_0}\prod_i x(Q_i)\equiv2^j\rho^4=2^j\ne0
 \quad\text{in }\mathbb F_3.
\]
In particular $v_3(C_0)=-E_0$, so $B_0$ is $3$-integral exactly when
all four indices are odd. This also proves $C_0\ne0$.

At $2$, equation~\eqref{further:eq:multipledepth} gives
$v(x(nP))=-3$ for odd $n$. If all four indices are odd, their product
has valuation $-12$; equation~\eqref{further:eq:componentvaluation}
then gives $v_2(C_0)\ge-12$. Hence $B_0$ is $2$-integral.
For every nonzero multiple $R$ of $P$, $v(f(R))>0$, so its two rational
coefficients are $2$-integral as well.

Replace the old physical anchor coordinate by
\begin{equation}\label{further:eq:newtheta}
 \Theta=q_2(f_0,f_1)+3\mathsf S(B_0^2+X^2),
\end{equation}
and retain the conic
\begin{equation}\label{further:eq:anchorconic}
 u^2+6v^2-2w^2-3\Theta^2+2\mathsf S^2=0.
\end{equation}
The source's quaternion criterion, used in
\nolinkurl{lem:v5-three-witness-anchor}, says that this conic is soluble
exactly when
\begin{equation}\label{further:eq:anchortest}
 q_2(f_0/C_{\rm a},f_1/C_{\rm a})+3(B_0^2+X^2)
 \in\mathbb Z_{(2)}\cap\mathbb Z_{(3)}.
\end{equation}

Both $q_2(u,v)$ and $u^2+v^2$ are anisotropic modulo $3$. Their
$3$-adic valuations are therefore twice the minimum valuation of their
arguments. The two summands in~\eqref{further:eq:anchortest} have,
respectively, even and odd valuations and cannot cancel.
Integrality at $3$ is equivalent to
\[
 f_0/C_{\rm a},\ f_1/C_{\rm a},\ B_0,\ X
 \quad\text{all being }3\text{-integral}.
\]
This first recovers the old depth anchor and forces all four root
indices to be odd. Only then use the $2$-adic coefficient bound above:
it makes $B_0$ $2$-integral, while $f_0/C_{\rm a}$ and
$f_1/C_{\rm a}$ are already $2$-integral because $C_{\rm a}$ is odd.
Thus integrality at $2$ in~\eqref{further:eq:anchortest} is equivalent
to $X$ being $2$-integral. A negative valuation of $3X^2$ cannot cancel
against the remaining integral terms.

The new anchor is consequently equivalent to the old depth condition,
the four odd-index conditions, and
$X\in\mathbb Z_{(2)}\cap\mathbb Z_{(3)}$.
This conclusion uses neither the global height argument nor the assertion
that $X$ is an integer. The original lower bound on $|a-1|$ is preserved.

\subsection{A fixed scalar correction}

Use the new label and correction
\begin{equation}\label{further:eq:newbeta}
 \beta=1224X^2+2451,
 \qquad 1224=2^3\cdot3^2\cdot17,
 \qquad c=21=3\cdot7.
\end{equation}
Remove the unused factor $2\beta-1$ from the scalar contact polynomial
and use
\begin{equation}\label{further:eq:r4}
 r=\beta(\beta-1)(\beta-2401)(32\beta+343).
\end{equation}
The justification that these four factors detect every required contact
is given below.
For every rational $X$, $\beta\ge2451>2401$, so $r>0$.
For integral $X$, $\beta$ is positive and odd and satisfies
\[
 \beta\equiv3\pmod8,\qquad
 \beta\equiv3\pmod9,\qquad
 \beta\equiv3\pmod{17}.
\]
These congruences make the local squareclasses of $r$ fixed.
At $2$, both $\beta-1$ and $\beta-2401$ have valuation one, and
\[
 \frac{\beta-1}{2}-\frac{\beta-2401}{2}=1200\equiv0\pmod8.
\]
Their odd parts therefore have product $1$ modulo $8$. The other two
factors have residues $3,7$, so $v_2(r)=2$ and $r/4\equiv5\pmod8$.
At $3$, only $\beta$ is divisible by $3$; it has valuation one and
$\beta/3\equiv1\pmod3$. The other residues are $2,2,1$, giving
$v_3(r)=1$ and $r/3\equiv1\pmod3$.
At $17$, the four residues are $3,2,16,14$, whose product is $1$.
Together with $21\equiv5\pmod8$, $21/3\equiv1\pmod3$, and
$21\equiv4\pmod{17}$, this proves
\[
 21r>0,\qquad
 21r\in\mathbb Q_2^{\times2}
      \cap\mathbb Q_3^{\times2}
      \cap\mathbb Q_{17}^{\times2}.
\]
For the last assertion, $21r\equiv4=2^2\pmod{17}$.
All assertions are for every integral $X$, not merely tested examples.

\paragraph{Input integrality.}
If the global argument gives $\beta\in\mathbb Z$ and $X=u/v$ in
lowest terms, then $v^2\mid1224$. Hence only $2$ or $3$ can divide $v$;
the exponent of $17$ in $1224$ is one. The new anchor has already made
$X$ integral at $2$ and $3$, so $v=1$.
Conversely, integral $X$ gives integral positive odd $\beta$.

\paragraph{The exact contact set and the removed factor.}
The source definition preceding \nolinkurl{hyp:o8-height} is
\[
 S_* =\{p:p\le76176\},\qquad
 \mathcal D=\{0,1,2401,-343/32\}.
\]
The radical $M_5(\beta)$ comprises the primes outside $S_*$ at which
$(\beta:1)$ has positive odd contact order with one of
$\infty,0,1,2401,-343/32$. The point $1/2$ is not a contact in this
definition.

For $p>76176$, the coefficient $1224$ is a unit. If $v_p(\beta)<0$,
then $v_p(\beta)=2v_p(X)$ is even. Thus the contact order with infinity,
$-v_p(\beta)$, is even, and no prime of $M_5(\beta)$ comes from infinity.
If $\beta$ is $p$-integral, a positive finite contact is
$v_p(\beta-\delta)>0$ for $\delta\in\mathcal D$.
The six differences between the finite contacts, up to sign, are
\[
 1,\quad 2401,\quad \frac{343}{32},\quad
 2400,\quad\frac{375}{32},\quad\frac{77175}{32}.
\]
Their numerator and denominator prime factors lie in $\{2,3,5,7\}$;
in particular $77175=3^2 5^2 7^3$.
Consequently at every consumed prime at most one finite contact has
positive valuation, and $32$ is a unit. Each odd finite contact makes
exactly one factor of~\eqref{further:eq:r4} have positive odd valuation.
In fact, for $p>76176$,
\[
 p\mid M_5(\beta)
 \quad\Longleftrightarrow\quad v_p(r)\text{ is odd},
\]
and in this case $v_p(r)>0$.

The discarded factor $2\beta-1$ is a unit at every such required
contact prime. Its values at the finite contacts are
$-1,1,4801,-359/16$; their numerator and denominator prime factors are
all at most $4801<76176$. Hence its removal does not even change the
valuation of $r$ at any consumed contact. It removes the extra tests
at $1/2$, which the height radical never required. No collision prime
outside the fixed cutoff is omitted.

\paragraph{Scalar norm soundness with four factors.}
Let $p\mid M_5(\beta)$ and $R_p=v_p(r)$, a positive odd integer.
The fixed correction $21$ is a unit at $p$. The scalar pole obstruction
applied to $b_{\rm sc}=21rW_{\rm sc}^2\ne0$ makes
$R_p+2v_p(W_{\rm sc})$ positive, so
$v_p(W_{\rm sc})>-R_p/2$.
In each norm node $N_r(u,v)=u^2+rv^2$, the two summands have opposite
valuation parity and cannot cancel. For $t<1$,
\[
 v_p(N_r(u,v))>-tR_p
 \ \Longrightarrow\ 
 v_p(u)>-tR_p/2,\quad v_p(v)>-(t+1)R_p/2.
\]
Starting with $t=1/2$ and descending the finite tree makes every
normalized leaf $l_i$ have valuation greater than $-R_p$.
For its unchanged order $m_i\in\{1,32\}$ and coefficient value $Y_i$,
$l_i=Y_i/r^{m_i}$, whence
\[
 v_p(Y_i)>(m_i-1)R_p\ge m_i-1,
 \qquad v_p(Y_i)\ge m_i.
\]
Zero inputs can be assigned valuation infinity and do not affect this
argument. Thus the same four coefficient tests put the shared and
scalar points in the formal group, and the same two order-$32$ tests
apply to their parameter ratio. The formal logarithm isometry and
$[h]R=[a-1]H$ give $v_p(a-1)\ge32$, exactly as in
\nolinkurl{v5:scalar32-bridge}. This proves
$M_5(\beta)^{32}\mid a-1$. The uniform theorem
\nolinkurl{thm:quartic-height19} then gives the unchanged bound
\[
 h(\beta)\le H_{19}|a-1|^{19/32}.
\]
No assertion about even negative valuations at $17$ is needed: that prime
is already excluded from the consumed-prime set.

\paragraph{Scalar completeness with four factors.}
For integral $X$, $r>1$ is an integer and $21r$ has the three local
squareclasses just proved. Use the finite prime support $\mathcal T$ of
$21r$. Choose the shared point deeply at all places above $\mathcal T$;
then choose $R=[a-1]P$ still deeper, with any additional prescribed
divisibility. The normalized leaves can be made integral when $p\mid r$
and divisible by $p$ when $p\mid21$ but $p\nmid r$.
The constant leaf $1$ is on the all-left path. Therefore the scalar
norm root $W_{\rm sc}$ is congruent to $1$ modulo every prime in
$\mathcal T$, as in the original scalar proof. It is a unit there.
The pole input
\[
 b_{\rm sc}=21rW_{\rm sc}^2
\]
has nonnegative valuation on $\mathcal T$ and even valuation outside it.
Its reduced denominator is square; it is positive and locally square
at $2,3,17$. The source's
\nolinkurl{lem:original-seven-parity-completeness} supplies its pole
witnesses under the same inherited restricted-converse and descent
hypotheses. That lemma requires local squares only at $2,3,17$, not
at every prime below $17$.

After $R$ is fixed, the same rational-function continuity argument in
\nolinkurl{v5:scalar32-bridge} supplies a progression for the shared index
preserving these finitely many valuations. Replacing five contact factors
by four introduces no new pole or forbidden point. The norm orders,
physical anchor, and subsequent comparison modulus are unchanged.

\paragraph{Where the global argument is reused.}
The local facts above establish the needed hypotheses before the global
argument starts. The packed anchor supplies odd roots and the original
scalar-index lower bound. The source's product-height result
\nolinkurl{v9:cor:scales}, local comparisons in
\nolinkurl{subsec:v9-local-depth}, Lucas step using
\nolinkurl{lem:o8-lucas}, and the two contradictions in\\
\nolinkurl{subsec:v9-soundness} use $\beta$ only as a rational label with
the stated scalar height bound. They therefore still force $n=\beta$.
The preceding denominator argument then forces $X\in\mathbb Z$.

Conversely, for integral $X$, take $n=\beta$, use the fixed correction
$21$, and choose the scalar index $s$ sufficiently divisible. Put
$a=s+1$ and choose the same positive odd Lucas labels as in
\nolinkurl{subsec:v9-completeness}. The conditions $a\equiv1\pmod3$
and $k\equiv1\pmod3$ in that construction concern $a$, not $n$.
All four roots are positive odd multiples, the old depth condition holds,
and integral $X$ makes the new anchor available. The subsequent fixed
modulus, prime-pattern, progression, continuity, and quotient-domain
arguments are unchanged. This is a transfer of the inherited global
proof, not a direct coordinate substitution into the old input label.

\subsection{One fewer coordinate for scalar preparation}

Let $b=b_{\rm sc}$ be the scalar pole input. With the new label and fixed
correction it is positive for every rational input satisfying the
arithmetic definitions: $r>0$, every node
$N_r(a,b)=a^2+rb^2$ is nonnegative, and the constant leaf $1$ on the
all-left path makes $W_{\rm sc}>0$. Thus $b\ne0$ holds independently
of any preparation witness.

The current preparation has five coordinates
\[
 t,\qquad t_2=t^2,\qquad T=(1+bt_2)/b,\qquad
 A=2/T,\qquad U_2=(A-b)^2.
\]
Its other pole equations use $U_2$ only. Replace these by four coordinates
$u,v,u_2,v_2$ satisfying
\begin{equation}\label{further:eq:preparationconic}
 u_2=u^2,\qquad v_2=v^2,\qquad
 u_2+bv_2-b^2=0.
\end{equation}
Use $u_2$ in the unchanged expression $e(17b-u_2)$.
Every residual in~\eqref{further:eq:preparationconic} is quadratic.

Given old witnesses, set
\[
 u=\frac{b(1-bt^2)}{1+bt^2},\qquad
 v=\frac{2bt}{1+bt^2}.
\]
The identity $(1-bt^2)^2+4bt^2=(1+bt^2)^2$ gives
$u^2+bv^2=b^2$ and preserves $u_2=U_2$.
Conversely, if $u+b\ne0$, set $t=v/(u+b)$. The conic gives
\[
 1+bt^2=\frac{2b}{b+u}\ne0,
\]
so the old preparation recovers $u$ and the same squared coordinate.
If $u=-b$, then $v=0$ because $b\ne0$. Replace $u$ by $+b$ and use
$t=0$, $T=1/b$, $A=2b$. The squared coordinate and every downstream
constraint are unchanged. This treats the one point omitted by the
finite rational parametrization and proves the two-way equivalence.

\subsection{Quadratic implementation and variable accounting}

The existing coordinate $X_2=X^2$ is reused in
$\beta=1224X_2+2451$ and in the additional anchor term
$3\mathsf S X_2$. No new product coordinate is required.
The physical anchor coordinate remains materialized; keep $X_2$ and the
existing coordinate $C_0^2$ materialized when lifting $\mathsf S$ to a
variable. Their products with $\mathsf S$ are then quadratic.
The anchor conic and its translation that cancels $\mathsf S^2$ retain
their formulas.

The arithmetic changes have the following counts, verified against the
final generated quadratic system:
\[
 \begin{array}{l|r}
 \text{Change}&\text{Coordinates removed}\\\hline
 \text{Six field guards replaced by rational guards}&6\\
 \text{Eight correction choices replaced by }21&8\\
 \text{Correction-times-}r\text{ gate becomes linear}&1\\
 \text{Conic scalar preparation}&1\\
 \text{Removed factor }2\beta-1&1\\\hline
 \text{Total}&17
 \end{array}
\]
These counts concern the additional arithmetic transformations only.
The final polynomial, coefficient-lifting count, and exact coefficient
certificates are separate construction outputs.


Combined with the separately verified sixteen-coordinate chart reduction,
these changes give
\[
 483-17-16=450
\]
total variables, including the input. Thus the new quartic has $449$
existential witnesses. The chart reduction and the final coefficient
certificates are checked separately in the construction package.
